% ----------------------------------------------------------------------------
% Experiment-wide settings. Edit here; every section reads these macros.
% ----------------------------------------------------------------------------
% Titles of the two reports (informe_aplicacion.tex and informe_ensayos.tex).
\newcommand{\TituloAplicacion}{VidFetch: seguimiento de robots en video}
\newcommand{\SubtituloAplicacion}{Uso, funcionamiento y código de la aplicación}
\newcommand{\TituloEnsayos}{Robots autopropulsados confinados: resultados de los ensayos filmados}
\newcommand{\SubtituloEnsayos}{Posición, velocidad, rotación y atascos en tres ensayos con 22 robots}
\newcommand{\Autores}{Dylan Frigerio}
\newcommand{\Institucion}{ITBA}
\newcommand{\FechaInforme}{\today}

% Real robot diameter (same value as "Diámetro real" in the app; 33 mm, consistent with the image
% measured with the enclosure scale). It sets the wall-contact radius and the area fraction, not the scale.
\newcommand{\DiametroRobotCM}{3.3}

% Enclosure (recinto): inner and outer diameter of the wall ring. The px -> mm scale of every frame
% comes from the inner diameter (same values as "Diámetro interior/exterior" in the app).
\newcommand{\DiametroIntRecintoCM}{18.5}
\newcommand{\DiametroExtRecintoCM}{19.5}

% Time scale of the time-lapse videos: frames that correspond to 1 s of real time
% (same value as "Escala de tiempo" in the app; 3 by default).
\newcommand{\CuadrosPorSegundo}{3}
